\documentclass[
    language=english,
    doctype=lecture-notes,
    institution=none,
]{../../omnilatex}

\RequirePackage{config/document-settings}

\begin{document}
\maketitle
\tableofcontents
\newpage

% Section 1: Introduction with a definition and theorem
\section{Introduction to Graph Theory}

\begin{definition}
A \textbf{graph} $G = (V, E)$ consists of a finite set $V$ of vertices
and a set $E \subseteq V \times V$ of edges.
\end{definition}

\begin{example}
Consider the graph $G = (\{1,2,3,4\}, \{(1,2),(2,3),(3,4),(4,1)\})$.
This is a simple cycle on four vertices, denoted $C_4$.
\end{example}

\begin{theorem}[Handshaking Lemma]
For any undirected graph $G = (V, E)$:
\[
    \sum_{v \in V} \deg(v) = 2|E|.
\]
\end{theorem}

\begin{proof}
Each edge contributes exactly 1 to the degree of each of its two
endpoints, so the total sum of degrees counts every edge twice.
\end{proof}

% Section 2: More content with lemmas and corollaries
\section{Trees and Connectivity}

\begin{definition}
A \textbf{tree} is a connected acyclic graph.  A \textbf{forest} is a
disjoint union of trees.
\end{definition}

\begin{lemma}
Every tree with $n$ vertices has exactly $n - 1$ edges.
\end{lemma}

\begin{proof}
By induction on $n$.  The base case $n = 1$ is trivial.
For the inductive step, remove a leaf vertex (which always exists
in a tree with $n \geq 2$).  The resulting graph is a tree on
$n - 1$ vertices with $(n-1) - 1 = n - 2$ edges by the inductive
hypothesis.  Adding back the leaf restores one edge, giving $n - 1$.
\end{proof}

\begin{corollary}
A connected graph with $n$ vertices and $n - 1$ edges is a tree.
\end{corollary}

\begin{theorem}[Euler's Formula for Planar Graphs]
For any connected planar graph with $V$ vertices, $E$ edges, and $F$ faces:
\[
    V - E + F = 2.
\]
\end{theorem}

\begin{proof}[Proof sketch]
Start with a spanning tree ($E = V - 1$, $F = 1$).  Each additional
edge adds exactly one face.  By induction, $V - (V-1+k) + (1+k) = 2$.
\end{proof}

% Section 3: Worked examples
\section{Worked Examples}

\begin{example}
Determine whether the graph $K_{3,3}$ (the complete bipartite graph
with three vertices in each partition) is planar.

\textbf{Solution.}  $K_{3,3}$ has $V = 6$ vertices and $E = 9$ edges.
If planar, by Euler's formula: $F = 2 - V + E = 2 - 6 + 9 = 5$.
But every face in a bipartite graph has at least 4 edges, so
$2E \geq 4F \implies 18 \geq 20$, a contradiction.  Hence $K_{3,3}$
is \textbf{not} planar.
\end{example}

\begin{remark}
The same technique shows that $K_5$ is not planar: $V = 5$, $E = 10$,
$F = 7$, and $2E = 20 \geq 3F = 21$, contradiction.
\end{remark}

\end{document}
